\documentclass[10pt,UTF8]{article}

\usepackage{geometry}


\usepackage{ctex}

\geometry{a4paper,scale=0.9}
\ctexset{today=old}

\usepackage[bookmarksnumbered=true]{hyperref}  % 生成大纲


\usepackage{times}

% \usepackage{graphicx} %插入图片的宏包
% \usepackage{subfigure}
% \usepackage{float} %设置图片浮动位置的宏包

\usepackage{amsmath}
\usepackage{amssymb}
\usepackage{mathtools}

\usepackage{physics}

% \usepackage{siunitx}
% % 关闭警告
% \ExplSyntaxOn
% \msg_redirect_name:nnn { siunitx } { physics-pkg } { none }
% \ExplSyntaxOff

%%%%%%%%%%%%%%%%%%%%%%%% tcolorbox %%%%%%%%%%%%%%%%%%%%%%%%%%
\usepackage[most]{tcolorbox}

\newenvironment{tipbox}
{\begin{tcolorbox}[colback=green!5!white,colframe=green!75!black]}
{\end{tcolorbox}}

\newenvironment{conceptbox}
{\begin{tcolorbox}[colback=blue!5!white,colframe=blue!75!black]}
{\end{tcolorbox}}

\newenvironment{thmbox}
{\begin{tcolorbox}[colback=yellow!5!white,colframe=yellow!75!black]}
{\end{tcolorbox}}

\newenvironment{definition box}
{\begin{tcolorbox}[colback=blue!5!white,colframe=blue!75!black,parbox=false,before upper=\par,before lower=\par]}
{\end{tcolorbox}}

\newenvironment{theorem box}
{\begin{tcolorbox}[colback=yellow!5!white,colframe=yellow!75!black,parbox=false,before upper=\par,before lower=\par]}
{\end{tcolorbox}}
%%%%%%%%%%%%%%%%%%%%%%%%%%%%%%%%%%%%%%%%%%%%%%%%%%%%%%%%%%%%
            
%%%%%%%%%%%%%%%%%%%%% amsthm %%%%%%%%%%%%%%%%
\usepackage{amsthm}

\theoremstyle{definition}
\newtheorem{definition inner}{Definition}[section]
\newenvironment{definition}[1][]
{\begin{definition box}\begin{definition inner}[#1]\hfill\par}
{\end{definition inner}\end{definition box}}

\theoremstyle{plain}
\newtheorem{theorem inner}{Theorem}[section]
\newenvironment{theorem}[1][]
{\begin{theorem box}\begin{theorem inner}[#1]\hfill\par}
{\end{theorem inner}\end{theorem box}}

\theoremstyle{definition}
\newtheorem*{proof inner}{\textit{Proof}}
\renewenvironment{proof}[1][]
{\begin{proof inner}[#1]\hfill\par}
{\qed\end{proof inner}}

\theoremstyle{remark}
\newtheorem*{remark}{\textit{Remark}}
%%%%%%%%%%%%%%%%%%%%%%%%%%%%%%%%%%%%%%%%%%%%%%


\newcommand{\mycases}[2][0.8]{
    \left\{\vphantom{\scalebox{#1}{
        $\begin{matrix}#2\end{matrix}$
    }}\right.\!\!\!\!
    \begin{array}{l@{\quad}l}#2\end{array}
}




\begin{document}

\section{Basic Concepts}

\paragraph{Experiment}
An \textbf{experiment} is a procedure that can be \emph{infinitely repeated} and has a well-defined set of \textbf{outcomes}.

\paragraph{Outcome}
An \textbf{outcome} is a possible result of an experiment.

\paragraph{Random Experiment}
A \textbf{random experiment} is an experiment that has more than one possible outcome.

\paragraph{Sample Space}
The set of all possible outcomes of an experiment is called the \textbf{sample space} of the experiment.

\paragraph{Event}
An \textbf{event} is a subset of the sample space.

When the random experiment is performed and the outcome is in a certain event,
we say that this event \emph{has occurred}.


\section{An Intuitive Definition of Probability}

\begin{enumerate}
    \item Repeat the experiment $n$ times.
    \item Count the number of times the event $A$ occurs and denote it by $\mathcal{N}(A)$.
\end{enumerate}

Then $\dfrac{\mathcal{N}(A)}{n}$ is called the \emph{relative frequency} of $A$ in the $n$ repetitions, and the \emph{probability} of $A$ is defined as
\[
P(A) := \lim_{n\to\infty}\frac{\mathcal{N}(A)}{n}
\]


\section{Set Theory}

\subsection{Some Notations}

\emph{Events} and \emph{sets} are equivalent in this context.
Thus usually we consider the \emph{sample space} $S$ to be the \emph{universe}
and have the following notations:

\begin{itemize}
    \item The \textbf{null} or \textbf{empty} set: $\emptyset$
    \item $A$ is a \textbf{subset} of $B$: $A \subseteq B$
    \item The \textbf{union} of $A$ and $B$: $A \cup B$
    \item The \textbf{intersection} of $A$ and $B$: $A \cap B$
    \item The \textbf{complement} of $A$ in $S$: $A'$
\end{itemize}

\begin{tipbox}
Note that the \emph{relative complement} of $A$ in $B$ can be obtained by
\[B\setminus A = B \cap A'\]
\end{tipbox}

\subsection{Some Concepts}

\subsubsection{Mutually Exclusive Events}
\begin{definition}[Mutually Exclusive Events]
    $A_1, A_2, \cdots, A_k$ are \textbf{mutually exclusive} $\iff$
    They are \textbf{disjoint} sets $\iff$
    $A_i \cap A_j = \emptyset$ for all $i \neq j$.
\end{definition}

\subsubsection{Exhaustive Events}
\begin{definition}[Exhaustive Events]
    $A_1, A_2, \cdots, A_k$ are \textbf{exhaustive} $\iff$
    $A_1 \cup A_2 \cup \cdots \cup A_k = S$.
\end{definition}

\subsubsection{Countable Sets}
\begin{definition}[Countable]
    A set is \textbf{countable} if there exists an \emph{injection} from it into $\mathbb{N}$.
\end{definition}

\begin{remark}
    Note that a countable set can be either \emph{finite} or \emph{infinite}.
\end{remark}

\subsection{Some Laws}

\subsubsection{Distributive Laws}
\begin{center}
    $A \cup (B \cap C) = (A \cup B) \cap (A \cup C)$ \\
    $A \cap (B \cup C) = (A \cap B) \cup (A \cap C)$
\end{center}

\begin{theorem}[Distributive Laws]
    \begin{center}
        $A \cup \left(B_1 \cup B_2 \cup \cdots \right) = 
        \left(A \cup B_1\right) \cup \left(A \cup B_2\right) \cup \cdots$ \\
        $A \cap \left(B_1 \cap B_2 \cap \cdots \right) = 
        \left(A \cap B_1\right) \cap \left(A \cap B_2\right) \cap \cdots$
    \end{center}
\end{theorem}

\subsubsection{De Morgan's Laws}
\begin{center}
    $(A \cup B)' = A' \cap B'$ \\
    $(A \cap B)' = A' \cup B'$
\end{center}

\begin{theorem}[De Morgan's Laws]
    \begin{center}
        $\left(A_1 \cup A_2 \cup \cdots \right)' = A_1' \cap A_2' \cap \cdots$ \\
        $\left(A_1 \cap A_2 \cap \cdots \right)' = A_1' \cup A_2' \cup \cdots$
    \end{center}
\end{theorem}

\subsubsection{Inclusion-Exclusion Principle}
\begin{theorem}[Inclusion-Exclusion Principle]
    Let $A_1, A_2, \cdots, A_n$ be finite sets and
    $\abs{A_i}$ denotes the cardinality of the set $A_i$, then
    \begin{gather*}
        \abs{\bigcup_{i=1}^n A_i} = 
        \sum_{\substack{J \subseteq \{1,2,\cdots,n\}\\J\neq\emptyset}}(-1)^{\abs{J}+1}
        \abs{\bigcap_{j\in J} A_j}
    \end{gather*}
\end{theorem}

\section{Probability and Probability Function}

\subsection{Definition}

\begin{definition}[Probability]
    \textbf{Probability} is a \emph{real-valued set function} that assigns
    a number $P(A)$ to each event $A$ in the sample space $S$.
    $P(A)$ is called the \textbf{probability} of $A$,
    and satisfies the following 3 properties (\emph{Probability Axioms}):
    
    \begin{itemize}
        \item $P(A) \geqslant 0$
        \item $P(S) = 1$
        \item For a \emph{countable} number of \emph{mutually exclusive} events
        $A_1, A_2, A_3, \cdots$,
        \[P(A_1\cup A_2\cup A_3\cup \cdots)=P(A_1)+P(A_2)+P(A_3)+\cdots\]
    \end{itemize}
\end{definition}
    
\subsection{Some Properties}

\subsubsection{Inclusion-Exclusion Principle}

\begin{theorem}[Inclusion-Exclusion Principle]
    If $A_1, A_2, \cdots, A_n$ are $n$ events, then
    \begin{gather*}
        P\left(\bigcup_{i=1}^n A_i\right) = 
        \sum_{\substack{J \subseteq \{1,2,\cdots,n\}\\J\neq\emptyset}}(-1)^{\abs{J}+1}
        P\left(\bigcap_{j\in J} A_j\right)
    \end{gather*}
\end{theorem}
    
\subsubsection{A Frequently Used Trick}

\begin{theorem}
    \begin{center}
        $P(A) = P(A \cap B) + P(A \cap B')$
    \end{center}
\end{theorem}

\subsection{Example}

Let $S=\left\{e_1, e_2, \cdots, e_m\right\}$ be a sample space,
where each $e_i$ is an outcome of the experiment.
Let $N(A)$ denote the number of outcomes in the event $A$.
It is easy to prove that the set function $P(A)$ defined as follows is a probability function:
\begin{itemize}
    \item $P(\left\{e_i\right\}) = \dfrac{1}{m}$ for all $i$
    \item $P(A) = \dfrac{N(A)}{N(S)}$ for any event $A$
\end{itemize}



\section{Probability Space}

\begin{definition}[Probability Space]
    \noindent
    A \textbf{probability space} is a triple $(S, \mathcal{F}, P)$ that consists of:
    \begin{itemize}
        \item A \textbf{sample space} $S$, which is a set of all possible outcomes.
        \item An \textbf{event space} $\mathcal{F}$, which is a set of all possible events.
        \item A \textbf{probability function} $P: \mathcal{F} \to [0,1]$ that satisfies the \emph{Probability Axioms}.
    \end{itemize}
\end{definition}





\end{document}
